\setpart{Pi0.5 · 实时方案}

\begin{frame}[t]{控制循环与推理循环分开运行}
\vspace{-.15cm}
\begin{tikzpicture}[x=1cm,y=.8cm,
  s/.style={box=Gold,font=\scriptsize,minimum width=2.0cm,minimum height=.62cm},
  p/.style={box=Teal,font=\scriptsize,minimum height=.66cm},
  ]
  % 传感器
  \node[s] (c) at (1.0,0) {相机 15 / 7.5 Hz};
  \node[s] (st) at (3.25,0) {本体状态 100 Hz};
  \node[s] (tk) at (5.5,0) {语言指令（事件）};
  % 观测槽
  \node[p,minimum width=7.6cm] (slot) at (3.25,-1.1) {最新观测槽：每路只保留最新一条，记录时间戳，读取时不等待};
  \draw[flow] (c) -- (c|-slot.north);
  \draw[flow] (st) -- (st|-slot.north);
  \draw[flow] (tk) -- (tk|-slot.north);
  % 第三行
  \node[p,minimum width=2.1cm] (trg) at (1.2,-2.4) {触发器\\{\tiny 计数 / 阈值 / 周期}};
  \node[p,minimum width=2.4cm,fill=Teal!14] (req) at (4.0,-2.4) {推理请求\\{\tiny 一次在途，带快照时刻}};
  \node[p,minimum width=2.3cm,fill=Orange!14,draw=Orange] (fw) at (6.7,-2.4) {框架推理\\{\tiny M* / vLLM / SGLang}};
  \node[p,minimum width=2.4cm] (aq) at (9.5,-2.4) {动作队列\\{\tiny 块 + 代数号，旧代数丢弃}};
  \draw[flow] (slot.south) -- (req.north);
  \draw[flow] (trg) -- (req);
  \draw[flow] (req) -- (fw);
  \draw[flow] (fw) -- (aq);
  % 第四行：控制
  \node[p,minimum width=2.4cm,fill=Gold!14,draw=Gold] (ct) at (9.5,-3.9) {控制时钟 15 Hz\\{\tiny 每 66.7 毫秒取一个动作}};
  \node[p,minimum width=2.4cm] (rb) at (1.2,-3.9) {机器人\\{\tiny 欠载时保持最后设定点（假设）}};
  \draw[flow] (aq) -- (ct);
  \draw[flow] (ct) -- (rb);
  \node[font=\tiny,text=Muted,anchor=west] at (3.0,-4.9) {推理慢只让动作变旧或欠载，不会让控制时钟停下};
\end{tikzpicture}
\par\takeaway{推理是异步的：控制时钟不等推理，推理也不等传感器全部到齐。两者只通过“最新观测槽”和“动作队列”交接。}
\par\memo{来源：本项目 src/\allowbreak{}runtime/\allowbreak{}observation.py、scheduler.py、action\_\allowbreak{}manager.py、controller.py；参考做法：LeRobot 观测槽与 ActionQueue（remote\_\allowbreak{}inference/\allowbreak{}engine.py:521--528；policies/\allowbreak{}rtc/\allowbreak{}action\_\allowbreak{}queue.py）。}
\end{frame}

\begin{frame}[t]{动作块怎么衔接：对齐、朴素与欠载}
\vspace{-.1cm}
\begin{tikzpicture}[x=1cm,y=.62cm]
  \def\X{1.75}
  \def\S{4.9}
  \def\dt{0.0667}
  % 时间刻度
  \foreach \tm in {0,0.5,1.0,1.5}{
    \draw[Line,line width=.4pt] ({\X+\tm*\S},-3.7) -- ({\X+\tm*\S},.65);
    \node[font=\tiny,text=Muted] at ({\X+\tm*\S},-3.95) {\tm\,s};
  }
  % 快照与到达
  \draw[Red,densely dashed,line width=.7pt] ({\X+.45*\S},.65) -- ({\X+.45*\S},-3.7);
  \node[font=\tiny,text=Red,anchor=south] at ({\X+.45*\S},.66) {快照 0.45 s};
  \draw[Blue,densely dashed,line width=.7pt] ({\X+.6*\S},.65) -- ({\X+.6*\S},-3.7);
  \node[font=\tiny,text=Blue,anchor=south west] at ({\X+.6*\S+.02},.66) {回复到达 0.6 s};
  % 车道 a：对齐执行
  \node[anchor=west,font=\scriptsize\bfseries,text=Ink] at (0,0) {对齐执行};
  \node[anchor=west,font=\tiny,text=Muted,align=left] at (0,-.42) {LeRobot aligned\\按当前时刻取动作};
  \foreach \i in {0,...,8}{\fill[Teal!60] ({\X+\i*\dt*\S},-.12) rectangle ({\X+(\i+1)*\dt*\S-.02},-.5);}
  \foreach \i in {9,...,14}{\fill[Teal!12] ({\X+\i*\dt*\S},-.12) rectangle ({\X+(\i+1)*\dt*\S-.02},-.5);}
  \foreach \i in {0,...,2}{\fill[Muted!30] ({\X+.45*\S+\i*\dt*\S},-.85) rectangle ({\X+.45*\S+(\i+1)*\dt*\S-.02},-1.23);}
  \foreach \i in {3,...,14}{\fill[Orange!70] ({\X+.45*\S+\i*\dt*\S},-.85) rectangle ({\X+.45*\S+(\i+1)*\dt*\S-.02},-1.23);}
  \node[anchor=west,font=\tiny,text=Muted] at ({\X+1.0*\S+.1},-1.04) {灰：跳过的动作};
  % 车道 b：朴素执行
  \node[anchor=west,font=\scriptsize\bfseries,text=Ink] at (0,-1.9) {朴素执行};
  \node[anchor=west,font=\tiny,text=Muted,align=left] at (0,-2.32) {OpenPI 参考客户端\\从动作 0 开始};
  \foreach \i in {0,...,14}{\fill[Teal!60] ({\X+.6*\S+\i*\dt*\S},-1.95) rectangle ({\X+.6*\S+(\i+1)*\dt*\S-.02},-2.33);}
  \node[anchor=west,font=\tiny,text=Ink] at ({\X+.6*\S+.5},-2.6) {每个动作比计划晚一个回复延迟（150 毫秒）};
  % 车道 c：回复晚到
  \node[anchor=west,font=\scriptsize\bfseries,text=Ink] at (0,-3.2) {回复晚到};
  \node[anchor=west,font=\tiny,text=Muted,align=left] at (0,-3.62) {欠载};
  \foreach \i in {0,...,14}{\fill[Teal!60] ({\X+\i*\dt*\S},-3.2) rectangle ({\X+(\i+1)*\dt*\S-.02},-3.58);}
  \fill[Red!70] ({\X+1.0*\S},-3.2) rectangle ({\X+1.3*\S},-3.58);
  \node[font=\tiny,text=Red] at ({\X+1.15*\S},-3.36) {欠载};
  \foreach \i in {13,14}{\fill[Orange!70] ({\X+1.3*\S+(\i-13)*\dt*\S},-3.2) rectangle ({\X+1.3*\S+(\i-12)*\dt*\S-.02},-3.58);}
\end{tikzpicture}
\par\takeaway{对齐执行避免了朴素执行的固定延迟；但回复一晚，剩下的动作很少，很快又欠载。示意图的时间为说明用，不是测量值。}
\par\memo{来源：设计文档 §2.7（朴素与对齐执行的预定时刻滞后：M* +68 毫秒，SGLang 分桶 +135 毫秒，vLLM-Omni +337 毫秒，对齐下约 0）；action\_\allowbreak{}manager.py 的 locate() 与 execute()。}
\end{frame}

\begin{frame}[t]{主流的几种做法（一）：何时发起推理}
\scriptsize
\renewcommand{\arraystretch}{1.12}
\setlength{\tabcolsep}{3pt}
\begin{tabularx}{\textwidth}{>{\bfseries}p{2.3cm} p{3.3cm} Y}
\toprule
\rowcolor{Mist}
做法 & 项目（版本） & 何时发起推理 \\
\midrule
计数触发 & OpenPI 参考客户端（215abfb） & 每执行固定步数就请求：DROID 满 8 步，LIBERO 每 5 步；与时间和剩余动作无关 \\
回放阈值 & LeRobot v0.7.0（410c095） & 队列剩余的播放时间 $\le$ 取 0.5 秒与（往返时间 + 一个周期）中较大者；pi05 块 50 步，30 fps 时剩 15 步请求 \\
异步 + RTC & openpi05-bsp（dfb7444）；rtc-pi05（HF，MIT） & 按时间：执行到 $t \ge \max(s_{\min}, d)$ 才发起，$d$ 取最近 10 次延迟的最大值；推理在后台线程进行 \\
多传感器对齐 & ROS 2 message\_filters（8.0.5） & 不触发推理；只把传感器消息按时间戳配成一组 \\
\bottomrule
\end{tabularx}
\par\takeaway{计数触发的参考客户端在推理期间阻塞控制循环；回放阈值与异步做法让推理在后台进行，控制循环不等推理。}
\par\memo{来源：agents\_\allowbreak{}results/2026-11-20\_\allowbreak{}openpi-remote-pi05.md 第 2 节；2026-11-20\_\allowbreak{}lerobot-async-rtc.md Q1（execution.py:228--230；factory.py:118）；2026-11-21\_\allowbreak{}level2-research-deep.md A、C 节。}
\end{frame}

\begin{frame}[t]{主流的几种做法（二）：新动作怎么接，能否中止}
\scriptsize
\renewcommand{\arraystretch}{1.12}
\setlength{\tabcolsep}{3pt}
\begin{tabularx}{\textwidth}{>{\bfseries}p{2.3cm} Y Y p{3.0cm}}
\toprule
\rowcolor{Mist}
做法 & 新块怎么接 & 能否中止或修正 & 代价 \\
\midrule
计数触发 & 从动作 0 开始，先进先出；无时间戳，无年龄检查 & 无中止，无 RTC & 每个动作晚一个推理延迟 \\
回放阈值 + 动作队列 & 观测槽只留最新一条；带代数号与游标的队列，对齐替换未执行部分 & 只在客户端丢弃；pi05 默认不开 RTC（需非 chunk 模式） & 回复慢时缓冲不够，动作断档 \\
异步 + RTC & 丢掉已执行的 $d$ 步，新块前缀按已发出的动作做推理期引导 & 服务端仍串行，无取消；openpi05-bsp 的 README 尚无数字 & 去噪每步求导，推理变慢 \\
训练期 RTC & 训练时采样延迟、前缀标为干净；推理时硬替换前缀 & 需要重新训练；EagleVLA 为未来修正模块，机制不同 & 推理开销低，但要训练好的检查点 \\
\bottomrule
\end{tabularx}
\par\takeaway{没有一个开源项目同时具备：服务端在推理中接收新观测、真正中止 GPU 推理，并带 RTC。RTC 都在客户端或模型里实现。}
\par\memo{来源：openpi 的 examples/droid、action\_\allowbreak{}chunk\_\allowbreak{}broker；LeRobot execution.py:204、539--582，configuration\_\allowbreak{}rtc.py:70；rtc-pi05 rtc/\allowbreak{}controller.py:146--183；Kinetix src/\allowbreak{}model.py:267--288。论文数字为作者自报，未复现。}
\end{frame}

\begin{frame}[t]{RTC：让新动作接上已经执行的动作}
\vspace{-.1cm}
\begin{minipage}[t]{5.6cm}
\minipagetop
\scalebox{.92}{\begin{tikzpicture}[x=1cm,y=.55cm]
  \def\X{0.2}
  \def\dd{.7}
  \draw[Line,line width=.4pt] (\X,-3.4) -- (\X+5.1,-3.4);
  \node[font=\tiny,text=Muted,anchor=north east] at (\X+5.1,-3.5) {动作序号};
  % 旧块：已执行（前 d 步）与未执行
  \foreach \i in {0,...,5}{\fill[Muted!45] ({\X+\i*.45},0) rectangle ({\X+\i*.45+.36},.5);}
  \node[font=\tiny,text=Muted,anchor=south] at ({\X+1.1},.55) {旧块已执行的部分};
  % 新块
  \foreach \i in {0,...,9}{\fill[Orange!65] ({\X+\i*.45+1.5},-1.2) rectangle ({\X+\i*.45+1.86},-.7);}
  \node[font=\tiny,text=Orange,anchor=north] at ({\X+3.0},-1.25) {新块（来自新的观测）};
  % 跳变
  \draw[Red,line width=.9pt,<->] ({\X+2.25},.5) -- ({\X+2.25},-.7);
  \node[font=\tiny,text=Red,anchor=west] at ({\X+2.32},-.1) {无 RTC：边界跳变};
  % 冻结的前缀（RTC）
  \fill[Teal!25] ({\X+1.5},-2.0) rectangle ({\X+2.7},-2.5);
  \node[font=\tiny,text=Teal,anchor=west] at ({\X+2.75},-2.25) {前缀 d 步：冻结或引导};
  \node[font=\tiny,text=Muted,anchor=west] at (\X,-2.85) {d = 推理耗时 ÷ 控制周期，向上取整};
\end{tikzpicture}}
\end{minipage}\hfill
\begin{minipage}[t]{5.0cm}
\minipagetop
\scriptsize
\renewcommand{\arraystretch}{1.1}
\begin{tabularx}{\linewidth}{>{\bfseries}p{1.7cm} Y}
\toprule
\rowcolor{Mist}
方式 & 做法与代价 \\
\midrule
无 RTC & 新块直接替换旧块；边界可能跳变 \\
推理期 guided & 前缀误差通过去噪步的 VJP 修正速度场，权重以 $\beta$ 截断（公式见下方）；不需要重训，但去噪步必须可求导 \\
训练期 trained & 训练时随机采样延迟，前缀标为干净动作；推理时硬替换；需要重训 \\
\bottomrule
\end{tabularx}
\end{minipage}
\par\begin{center}\scriptsize 权重 $w_t=\min\big(\frac{t^2+(1-t)^2}{t(1-t)},\,\beta\big)$\end{center}
\par\takeaway{guided 不需要重训，但我们的 M* 与 SGLang core 去噪步都在 CUDA 图里，求导需要重新捕获图；这是引擎级改动。}
\par\memo{来源：Kinetix src/\allowbreak{}model.py:236--258；rtc-pi05 rtc/\allowbreak{}masking.py:122--139；LeRobot configuration\_\allowbreak{}rtc.py:\allowbreak{}41--52，\allowbreak{}modeling\_\allowbreak{}rtc.py:\allowbreak{}46--48；\allowbreak{}设计文档 §0 第 7 条。}
\end{frame}
